\chapter{Теоријске основе}
\label{ch:teorija}

У овом поглављу дат је преглед кључних теоријских концепата на којима се темељи
истраживање. Најпре је објашњен појам реализоване волатилности као изведене, а не
посматране величине, заједно са начином на који се она конструише из интрадневних
приноса. Потом је представљен први експериментални фактор, дизајн мерења, кроз три
одлуке које га чине: избор фреквенције узорковања, избор естиматора отпорног на
скокове и избор теста за детекцију скокова, уз објашњење начина на који се непрекидна
и скоковита компонента раздвајају. Следи преглед изведених променљивих које из тих
одлука произлазе и чине непосредан улаз модела. Након тога изложен је други
експериментални фактор, архитектуре предиктивних модела, почев од економетријске HAR
породице и њених проширења, преко модела заснованих на стаблима, све до
трансформерских и темељних модела временских серија, са посебним освртом на
механизме којима сваки од њих обезбеђује ненегативну прогнозу. На крају су изложени
валидациони протокол и апарат за евалуацију, који обухвата функције губитка, тестове
статистичке значајности, економску проверу путем вредности при ризику и декомпозицију
варијансе као централни методолошки инструмент рада.

\section{Реализована волатилност као изведена величина}

\subsection{Полазни модел цене}

Претпоставља се да логаритам цене финансијског инструмента прати семимартингал са
скоковима, што је стандардна поставка у литератури о реализованој волатилности:

\begin{equation}
dp(t) = \mu(t)\,dt + \sigma(t)\,dW(t) + \kappa(t)\,dq(t)
\end{equation}

где је $p(t)$ логаритам цене, $\mu(t)$ компонента сноса, $\sigma(t)$ тренутна (спот)
волатилност, $W(t)$ стандардно Брауново кретање, $q(t)$ бројачки процес који
региструје скокове, а $\kappa(t)$ величина скока у тренутку $t$.

Величина од интереса јесте \textbf{квадратна варијација} процеса на дану $t$, која
се разлаже на непрекидни део (интегрисану варијансу) и део који потиче од скокова:

\begin{equation}
\QV_t = \int_{t-1}^{t}\sigma^2(s)\,ds + \sum_{t-1 < s \le t}\kappa^2(s)
\end{equation}

Прва сума представља \textbf{континуирану компоненту}, а друга \textbf{скоковиту
компоненту}. Централна чињеница за овај рад јесте да ниједна од ове две величине
није непосредно посматрана: обе се оцењују из дискретно узоркованих цена, а свака
оцена носи методолошке одлуке које се у наставку разматрају као експериментални
фактори.

\subsection{Реализована варијанса}

Ако се трговачки дан подели на $n$ интервала једнаке дужине, интрадневни
логаритамски принос на $i$-том интервалу дефинише се као

\begin{equation}
r_{t,i} = p_{t,i} - p_{t,i-1}
\end{equation}

\textbf{Реализована варијанса} јесте сума квадрата тих приноса:

\begin{equation}
\RV_t = \sum_{i=1}^{n} r_{t,i}^{2}
\end{equation}

Када $n \to \infty$ и у одсуству микроструктурног шума, $\RV_t$ конвергира ка
квадратној варијацији $\QV_t$, дакле ка збиру континуиране и скоковите компоненте.
У имплементацији се $\RV$ рачуна искључиво из приноса унутар једног трговачког дана:
прекноћни принос се никада не укључује, јер он не представља континуирано трговање
и његово укључивање би контаминирало меру догађајем који припада другој временској
скали.

\subsection{Реализована квартичност}

Реализована квартичност оцењује четврти момент и потребна је као мера прецизности
саме оцене $\RV$, због чега улази у HARQ модел (одељак~\ref{subsec:harq}):

\begin{equation}
\RQ_t = \frac{n}{3}\sum_{i=1}^{n} r_{t,i}^{4}
\end{equation}

Ова величина нестабилна је при малом $n$, што је разлог за апсолутни доњи праг броја
интрадневних опсервација уведен у контроли квалитета.

\subsection{Семиваријансе са предзнаком}

Патон и Шепард \parencite{pattonSheppard2015} показују да будућа волатилност зависи
различито од приноса различитог предзнака. Реализоване семиваријансе раздвајају
$\RV$ по предзнаку приноса:

\begin{equation}
\RS_t^{+} = \sum_{i=1}^{n} r_{t,i}^{2}\,\mathbf{1}\{r_{t,i} > 0\},
\qquad
\RS_t^{-} = \sum_{i=1}^{n} r_{t,i}^{2}\,\mathbf{1}\{r_{t,i} < 0\}
\end{equation}

Између њих важи тачан идентитет, који се у имплементацији проверава као тврда
провера исправности:

\begin{equation}
\RS_t^{+} + \RS_t^{-} = \RV_t
\label{eq:semivar-identity}
\end{equation}

\textbf{Варијација скока са предзнаком} дефинише се као њихова разлика:

\begin{equation}
\SJ_t = \RS_t^{+} - \RS_t^{-}
\end{equation}

Дневни принос, потребан за полужни (\textit{leverage}) члан и за евалуацију
вредности при ризику, јесте сума интрадневних приноса:

\begin{equation}
r_{d,t} = \sum_{i=1}^{n} r_{t,i}
\end{equation}

\section{Први експериментални фактор: Дизајн мерења}
\label{sec:dizajn}

Овај одељак описује прву половину факторског експеримента. Три одлуке које се у
литератури најчешће доносе прећутно овде се третирају као експлицитне
експерименталне димензије, означене као А1, А2 и А3.

\paragraph{Зашто су експериментална фактора два, а не четири.}
Рад разликује величине које се \textbf{варирају} од величина које из њих
\textbf{произлазе} или служе за \textbf{исецање} резултата. Дизајн мерења
(овај одељак) и архитектура модела (одељак~\ref{sec:modeli}) јесу праве
експерименталне осе: оне се укрштају у потпуни факторски план и управо се
њихови удели варијансе пореде у декомпозицији варијансе, што је централно
питање рада. Изведене променљиве (одељак~\ref{sec:izvedene}) нису
алтернативе једна другој, већ настају из дизајна мерења и заједно чине улаз
модела, па не представљају осу која се вари́ра. Стратификација
(одељак~\ref{subsec:stratifikacija}) уопште не утиче на генерисање прогноза,
него на начин на који се већ добијени резултати приказују, па припада
евалуацији. Овакво разграничење је неопходно да би удели варијансе у
декомпозицији уопште били интерпретабилни.

\subsection{Димензија А1 — фреквенција узорковања}

Избор мреже узорковања ставља истраживача пред компромис који нема решење важеће за
све случајеве.

Са статистичке стране, асимптотска варијанса оцене реализоване варијансе опада са
бројем опсервација:

\begin{equation}
\operatorname{Var}(\RV_t) \approx \frac{2}{n}\int_{t-1}^{t}\sigma^{4}(s)\,ds
\end{equation}

Ово сугерише да је гушћа мрежа увек боља. Међутим, посматрана цена није ефикасна
цена, већ садржи \textbf{микроструктурни шум}:

\begin{equation}
p_{t,i}^{obs} = p_{t,i}^{*} + u_{t,i}
\end{equation}

где је $p^{*}$ ефикасна (неопсервабилна) цена, а $u$ шум који потиче од распона
између куповне и продајне цене, дискретизације цена и ефеката извршења налога. Под
претпоставком да је шум независан и једнако распоређен, очекивана вредност
реализоване варијансе рачунате из посматраних цена износи

\begin{equation}
E\!\left[\RV_t^{obs}\right] = \QV_t + 2n\,E\!\left[u^{2}\right]
\label{eq:noise-bias}
\end{equation}

Други члан расте \textbf{линеарно са $n$}. Пораст фреквенције узорковања стога
смањује варијансу оцене, али истовремено линеарно увећава пристрасност. То је разлог
зашто је петоминутна мрежа постала стандард у литератури: она представља емпиријски
компромис, а не теоријски оптимум.

У овом раду мрежа узима вредности из скупа

\begin{equation}
A1 \in \{1m,\ 5m,\ 15m\}
\end{equation}

Свака мрежа изводи се преузорковањем истих учитаних једноминутних свећица, тако да
димензија А1 мења искључиво агрегацију, а не извор података. Очекивани број свећица по
дану приказан је у табели~\ref{tab:bars}.

\begin{table}[htbp]
\centering
\caption{Очекивани број свећица по мрежи и типу трговачког дана}
\label{tab:bars}
\begin{tabular}{lcc}
\toprule
Мрежа & Пун дан (390 мин.) & Скраћен дан (210 мин.) \\
\midrule
1m  & 390 & 210 \\
5m  & 78  & 42  \\
15m & 26  & 14  \\
\bottomrule
\end{tabular}
\end{table}

\subsection{Димензија А2 — естиматори отпорни на скокове}

Циљ сваког естиматора у овој групи јесте да оцени \textbf{искључиво континуирани
део} волатилности, дакле оно што би $\RV$ био да скокова нема. Разлика $\RV$ и
таквог естиматора представља оцену доприноса скокова.

\paragraph{Биповер варијација (BPV).}
Барндорф-Нилсен и Шепард \parencite{bns2004} уместо квадрата једног приноса множе
апсолутне вредности два суседна приноса. Пошто скок утиче на само два сабирка и то
без квадрирања, његов утицај је потиснут:

\begin{equation}
\BPV_t = \mu_1^{-2}\cdot\frac{n}{n-1}\sum_{i=2}^{n}\left|r_{t,i}\right|\left|r_{t,i-1}\right|
\end{equation}

где је константа скалирања

\begin{equation}
\mu_1 = E\!\left[|Z|\right] = \sqrt{\frac{2}{\pi}} \approx 0{,}7979,
\qquad
\mu_1^{-2} = \frac{\pi}{2} \approx 1{,}5708
\end{equation}

Фактор корекције за коначан узорак $n/(n-1)$ рачуна се из стварног броја приноса тог
дана након чишћења, никада из фиксиране константе: скраћени дани легитимно носе мањи
број свећица, па би фиксиран делилац уносио систематску пристрасност на сваки такав
дан.

\paragraph{Медијална реализована варијанса (MedRV).}
Андерсен, Добрев и Шаумбург \parencite{andersen2012} узимају медијану сваке тројке
узастопних апсолутних приноса, чиме се највећи члан тројке игнорише:

\begin{equation}
\MedRV_t = \frac{\pi}{6 - 4\sqrt{3} + \pi}\cdot\frac{n}{n-2}
\sum_{i=2}^{n-1}\med\!\left(|r_{t,i-1}|,\,|r_{t,i}|,\,|r_{t,i+1}|\right)^{2}
\end{equation}

Овај естиматор отпорнији је од BPV-а у ситуацији када се скокови јаве у узастопним
приносима, дакле у случају у коме је BPV-ов производ двају суседних приноса потпуно
контаминиран.

\paragraph{Биповер варијација са прагом (TrBPV).}
Трећа варијанта уклања приносе изнад прага изведеног из локалне волатилности пре
него што се формирају унакрсни производи:

\begin{equation}
\TrBPV_t = \mu_1^{-2}\cdot\frac{n}{n-1}
\sum_{i=2}^{n}\left|\tilde{r}_{t,i}\right|\left|\tilde{r}_{t,i-1}\right|
\end{equation}

где је принос са прагом дефинисан као

\begin{equation}
\tilde{r}_{t,i} = r_{t,i}\cdot\mathbf{1}\!\left\{|r_{t,i}| \le \lambda\hat{\sigma}_t\right\}
\end{equation}

Локална скала оцењује се робусно, преко медијане апсолутног приноса скалиране на
јединицу стандардне девијације под нормалношћу:

\begin{equation}
\hat{\sigma}_t = \frac{\med\left(|r_{t,i}|\right)}{0{,}6745}
\end{equation}

Праг је фиксиран на $\lambda = 3{,}0$ кроз цео простор дизајна. Ова одлука је
намерна: увођење прага као четврте димензије дизајна увећало би мрежу са 36 на 108
комбинација, а пошто се параметар односи само на један од три естиматора, дизајн би
постао неуравнотежен и удели варијансе по фактору у анализи варијансе постали би
неинтерпретабилни.

\subsection{Димензија А3 — тестови за детекцију скокова}
\label{subsec:a3}

\paragraph{Наивни праг.}
Дан се проглашава даном са скоком ако укупан дневни принос пређе фиксиран број
стандардних девијација:

\begin{equation}
S_t^{naive} = \frac{\left|r_{d,t}\right|}{\sqrt{\BPV_t}},
\qquad
\text{скок ако } S_t^{naive} > 3
\end{equation}

Кључан детаљ имплементације јесте да се стандардна девијација у имениоцу изводи из
естиматора отпорног на скокове, никада из стандардне девијације рачунате преко целог
дана. Друга варијанта укључила би сам скок који се тестира у сопствени именилац, па
би довољно велик скок надувао сопствени праг и избегао детекцију. Овај тест не
користи се као озбиљна метода, већ као доња граница квалитета.

\paragraph{БНС тест.}
Барндорф-Нилсен и Шепард \parencite{bns2006} пореде $\RV$ и естиматор отпоран на
скокове, скалирано триповер квартичношћу. Релативни допринос скокова дефинише се као

\begin{equation}
RJ_t = \frac{\RV_t - \BPV_t}{\RV_t}
\end{equation}

Триповер квартичност, потребна за нормализацију, износи

\begin{equation}
\TPQ_t = n\,\mu_{4/3}^{-3}\cdot\frac{n}{n-2}
\sum_{i=3}^{n}\left|r_{t,i}\right|^{4/3}\left|r_{t,i-1}\right|^{4/3}\left|r_{t,i-2}\right|^{4/3}
\end{equation}

Тест статистика има асимптотску стандардну нормалну расподелу под нултом хипотезом о
одсуству скокова:

\begin{equation}
Z_t = \frac{RJ_t}{\sqrt{\theta\cdot\dfrac{\TPQ_t}{n\,\BPV_t^{2}}}}\ \sim\ N(0,1)
\label{eq:bns}
\end{equation}

где је константа асимптотске варијансе

\begin{equation}
\theta = \left(\frac{\pi}{2}\right)^{2} + \pi - 5 \approx 0{,}6090
\end{equation}

Нормализација са $n\,\BPV_t^{2}$, а не са $\RV_t^{2}$, јесте оно што даје исправну
асимптотску скалу. Ово је провером потврђено: у ранијој верзији формуле, у којој је
фактор $n$ био изостављен, убачен скок величине двадесет стандардних девијација
производио је статистику од свега $0{,}69$, док иста провера након исправке даје
вредност око $8{,}0$.

\paragraph{Ли-Микландов тест.}
Ли и Микланд \parencite{leeMykland2008} не тестирају дан као целину, већ сваки
\textbf{интрадневни принос} посебно, делећи га локалном оценом волатилности
изведеном из клизног прозора од $K$ претходних опсервација:

\begin{equation}
L_{t,i} = \frac{r_{t,i}}{\hat{\sigma}_{t,i}}
\end{equation}

Локална оцена волатилности рачуна се биповер поступком, у сопственој конвенцији Лија
и Микланда, дакле \textbf{без} корекције фактором $\mu_1^{-2}$:

\begin{equation}
\hat{\sigma}_{t,i}^{2} = \frac{1}{K-1}\sum_{j=i-K+1}^{i-1}\left|r_j\right|\left|r_{j-1}\right|
\label{eq:local-sigma}
\end{equation}

Ова конвенција има важну последицу коју је неопходно навести. Пошто изостављена
корекција значи да $\hat{\sigma} \to \sigma\mu_1$, статистика под нултом хипотезом
није стандардно нормална, већ износи

\begin{equation}
L_{t,i} \to \frac{Z}{\mu_1} = Z\cdot 1{,}2533
\end{equation}

дакле она је за 25\% већа од $|Z|$. Критична вредност мора бити изведена уз ту
конвенцију. Максимум апсолутних статистика дана приближно прати Гумбелову расподелу,
са параметрима положаја и размере

\begin{equation}
C_n = \frac{\sqrt{2\ln n}}{c} - \frac{\ln\pi + \ln(\ln n)}{2c\sqrt{2\ln n}},
\qquad
S_n = \frac{1}{c\sqrt{2\ln n}}
\end{equation}

где је $c = \mu_1 = \sqrt{2/\pi}$. Критична вредност на нивоу поверења од 95\% износи

\begin{equation}
q_n = C_n + S_n\beta^{*},
\qquad
\beta^{*} = -\ln\!\left(-\ln 0{,}95\right)
\end{equation}

Дан се проглашава даном са скоком ако важи $\max_i |L_{t,i}| > q_n$. Критична
вредност зависи од $n$ и рачуна се за сваки дан посебно: фиксиран праг (на пример
$|L| > 4{,}0$) помешао би димензију А1 са димензијом А3, јер би исти праг значио различит
ниво значајности на различитим мрежама.

Дужина локалног прозора $K$ фиксирана је на \textbf{четири трговачке сесије на свакој
мрежи} ($K = 1560/312/104$ за мреже $1m/5m/15m$). Ово је намерна одлука, а не
преузимање вредности из литературе: када би се дужина прозора изражена у сесијама
мењала између мрежа, класификација скока зависила би од две ствари које се обе мењају
са димензијом А1, што би искривило удео варијансе приписан дизајну мерења.

\paragraph{Ли-Микланд са контролом стопе лажних открића.}
Гумбелова критична вредност већ представља контролу грешке на нивоу породице
\textbf{унутар једног дана}, преко свих $n$ интрадневних тестова тог дана. Проблем
који остаје јесте акумулација \textbf{преко дана}: и када је сваки дан појединачно
исправно контролисан, хиљаде дана заједно генеришу лажна открића.

Бенјамини-Хохбергова процедура \parencite{benjaminiHochberg1995} примењује се стога
преко клизног прозора од $W$ претходних дана плус текућег дана. За сортиране
$p$-вредности $p_{(1)} \le p_{(2)} \le \dots \le p_{(m)}$ проналази се највећи индекс
који задовољава

\begin{equation}
k^{*} = \max\left\{k\ :\ p_{(k)} \le \frac{k}{m}q\right\}
\end{equation}

а значајним се проглашавају сви тестови са $p \le p_{(k^{*})}$. Прозор износи
$W = 250$ дана, што одговара приближно једној трговачкој години. Клизни прозор је
обавезан: примена преко целог узорка класификовала би дан користећи $p$-вредности из
година које му тек следе, што представља гледање унапред.

Емпиријска провера (AAPL, мрежа 5m, период 2012--2013) показује да удео дана са
скоком пада са 31,9\% на 17,9\%, уз неслагање двају варијанти на 14,0\% дана.

\subsection{Условна декомпозиција на континуирану и скоковиту компоненту}
\label{subsec:decompose}

Ово је место на коме се димензије А2 и А3 срећу и једино место у ком се обе одлуке
истовремено уграђују у улазне податке модела.

\textbf{Безусловна} декомпозиција, каква се често среће, дефинисала би компоненте као
$J_t = \max(\RV_t - \BPV_t, 0)$ и $C_t = \RV_t - J_t$. Она зависи искључиво од
естиматора, а никада од теста: заменом БНС теста Ли-Микландовим, вредности $J$ и $C$
остале би идентичне, што би димензију А3 учинило потпуно безефектном за сваки улаз
модела. То би поништило саму сврху рада.

\textbf{Условна} декомпозиција \parencite{abd2007} то исправља тако што скоковиту
компоненту признаје само на данима на којима је тест значајан:

\begin{equation}
J_t = \mathbf{1}\{\text{тест значајан на дан } t\}\cdot\max\!\left(\RV_t - \widehat{\BPV}_t,\ 0\right)
\label{eq:cond-decomp}
\end{equation}

\begin{equation}
C_t = \RV_t - J_t
\end{equation}

Сада $J$ зависи и од тога који је естиматор употребљен (димензија А2) и од тога који је
тест прогласио дан даном са скоком (димензија А3). На данима без значајног теста важи
$J_t = 0$ и сва варијација приписује се континуираном делу, што је и теоријски
исправно: емпиријски је утврђено да просечан јаз $(\RV - \BPV)/\RV$ износи свега
3,5\%, што значи да је сирови јаз на данима без скока шум естиматора, а не скок.

Иста јединствена формула користи се за сваку од четири варијанте теста. Давање
другачије формуле једној варијанти (на пример сумирање квадрата означених приноса за
Ли-Микландов тест) помешало би „ефекат теста“ са „ефектом формуле“ у истом фактору
анализе варијансе.

\subsection{Простор дизајна мерења}

Комбиновањем трију фактора добија се простор дизајна

\begin{equation}
D = A1 \times A2 \times A3,
\qquad
|D| = 3 \times 3 \times 4 = 36
\end{equation}

Сваки дизајн идентификује се низом облика \texttt{мрежа\_\_естиматор\_\_тест} (на
пример \texttt{5m\_\_bpv\_\_bns}), који касније постаје кључ фактора у декомпозицији
варијансе.

\section{Изведене променљиве као улаз модела}
\label{sec:izvedene}

Величине описане у овом одељку не представљају експерименталну осу која се
вари́ра, него скуп променљивих које настају из дизајна мерења и заједно чине
улаз предиктивних модела. Оне се међусобно не искључују: недељни просек није
алтернатива дневној вредности, већ обе истовремено улазе у исту регресију.
Пошто свака од њих директно зависи од трију димензија из
одељка~\ref{sec:dizajn}, промена дизајна мерења мења њихове бројчане вредности,
а тиме и све што модели виде.

\subsection{HAR агрегати}

Из дневних мера граде се недељни и месечни просеци, који чине улаз HAR породице
модела:

\begin{equation}
\RV_{w,t} = \frac{1}{5}\sum_{j=0}^{4}\RV_{t-j},
\qquad
\RV_{m,t} = \frac{1}{22}\sum_{j=0}^{21}\RV_{t-j}
\end{equation}

Исти поступак примењује се и на компоненте $C$ и $J$, дајући величине
$C_d, C_w, C_m$ и $J_d, J_w, J_m$. Три временска хоризонта одговарају
Корсијевој хипотези о хетерогеном тржишту, према којој учесници реагују
различитим брзинама.

\subsection{Преглед скупа изведених променљивих}

Изведене променљиве приказане су у табели~\ref{tab:izvedene}, заједно са
назнаком од којих димензија дизајна мерења свака зависи. Та зависност је
суштинска за тумачење резултата: променљиве које зависе од димензија А2 и А3
мењају вредност са сваком изменом дизајна, док оне које зависе само од А1
остају идентичне кроз све комбинације естиматора и теста.

\begin{table}[htbp]
\centering
\caption{Изведене променљиве и њихова зависност од димензија дизајна мерења}
\label{tab:izvedene}
\begin{tabularx}{\textwidth}{L{3.4cm} X L{2.2cm}}
\toprule
Група & Променљиве & Зависи од \\
\midrule
HAR компоненте &
$\RV_d$, $\RV_w$, $\RV_m$ &
А1 \\
\addlinespace
Компоненте скокова &
$C_d$, $C_w$, $C_m$, $J_d$, $J_w$, $J_m$, бинарни индикатор скока &
А1, А2, А3 \\
\addlinespace
Семиваријансе &
$\RS^{+}$, $\RS^{-}$, варијација са предзнаком $\SJ$ &
А1 \\
\addlinespace
Мера прецизности &
реализована квартичност $\RQ$ &
А1 \\
\addlinespace
Полужни члан &
дневни принос $r_d$ &
А1 \\
\addlinespace
Егзогена променљива &
$\VIX$ са доцњом од једног дана &
--- \\
\bottomrule
\end{tabularx}
\end{table}

Основ за разлагање на семиваријансе јесте налаз Патона и Шепарда
\parencite{pattonSheppard2015} да је негативна семиваријанса знатно значајнија
за будућу волатилност од позитивне. Индекс имплицитне волатилности узима се
искључиво са доцњом од једног дана: коришћење вредности из истог дана за
предвиђање тог истог дана представљало би гледање унапред.

\section{Други експериментални фактор: Архитектуре предиктивних модела}
\label{sec:modeli}

\subsection{Конструкција циљне променљиве}
\label{subsec:target}

Циљна величина јесте \textbf{$h$-дневни просек} будуће реализоване варијансе, што је
Корсијева изворна HAR конструкција и оно што свака студија у овој области извештава:

\begin{equation}
y_{t,h} = \frac{1}{h}\sum_{i=1}^{h}\RV_{t+i}
\label{eq:target}
\end{equation}

Хоризонти су $h \in \{1, 5, 22\}$, што одговара дану, недељи и месецу. Циљ се моделује
у логаритамском простору:

\begin{equation}
\ln y_{t,h} = \ln\!\left(\frac{1}{h}\sum_{i=1}^{h}\RV_{t+i}\right)
\end{equation}

Разлог за логаритамску трансформацију јесте изражена десна асиметрија расподеле
реализоване волатилности. Мали број модела, чија је позитивност структурна, уместо
тога користи циљ у нивоима (одељци \ref{subsec:pozitivnost} и \ref{subsec:mem}).

\subsection{Проблем повратне трансформације и Јенсенова неједнакост}
\label{subsec:jensen}

Ако модел предвиђа $\ln y$, наивна повратна трансформација $\exp(\widehat{\ln y})$
\textbf{систематски потцењује} ниво. Разлог је Јенсенова неједнакост: за конвексну
функцију важи

\begin{equation}
E\!\left[\exp(X)\right] \ge \exp\!\left(E[X]\right)
\end{equation}

па је $\exp(\widehat{\ln y})$ оцена условне \textbf{медијане}, а не условне средине.
Пошто су QLIKE и RMSE минимизовани условном средином, потребна је корекција.

Уобичајена Гаусова корекција гласила би $a = \tfrac{1}{2}\sigma^{2}$, али она
претпоставља да су резидуали тачно нормално распоређени. Резидуали реализоване
волатилности то не морају бити, посебно код модела заснованих на стаблима. У овом
раду корекција је стога \textbf{непараметарска}:

\begin{equation}
a = \ln\!\left(\frac{1}{m}\sum_{k=1}^{m}\exp(e_k)\right),
\qquad
e_k = \ln y_k - \widehat{\ln y}_k
\label{eq:correction}
\end{equation}

Повратна трансформација гласи

\begin{equation}
\hat{y} = \exp\!\left(\widehat{\ln y} + a\right)
\end{equation}

Корекција $a$ оцењује се \textbf{искључиво на валидационом преклопу} (последња година
тренинга), никада на резидуалима унутар узорка: резидуали модела заснованих на
стаблима унутар узорка близу су нуле, што би такве моделе тихо недовољно кориговало у
односу на моделе оцењене методом најмањих квадрата.

За моделе који предвиђају \textbf{збир компоненти}, корекција мора бити оцењена на
већ сабраној прогнози, и то као мултипликативни фактор:

\begin{equation}
m^{*} = \frac{1}{N}\sum_{k=1}^{N}\frac{y_k}{\hat{y}_k^{comb}},
\qquad
\hat{y}^{corr} = m^{*}\hat{y}^{comb}
\end{equation}

Примена засебне корекције по компоненти и сабирање након корекције оставило би збир
без икакве контролисане корекције, јер Јенсенова неједнакост није адитивна преко
компоненти.

\subsection{Три различита механизма позитивности}
\label{subsec:pozitivnost}

Волатилност је по дефиницији ненегативна, па сваки модел мора гарантовати ненегативну
прогнозу. У овом раду постоје \textbf{три намерно различита механизма}, који се никада
не обједињују, управо зато што рад испитује да ли сам механизам позитивности, а не
само скуп улазних променљивих, утиче на тачност. У табели~\ref{tab:pozitivnost}
приказана су сва три механизма са припадајућим моделима.

\begin{table}[htbp]
\centering
\caption{Три механизма обезбеђивања ненегативне прогнозе}
\label{tab:pozitivnost}
\begin{tabularx}{\textwidth}{L{3.4cm} L{4.6cm} X}
\toprule
Механизам & Модели & Начин гарантовања \\
\midrule
Логаритам и непараметарска корекција &
\texttt{log\_har}, \texttt{char}, \texttt{har\_j}, \texttt{shar}, \texttt{harq},
\texttt{har\_iv}, \texttt{arfima}, стабла, PatchTST &
Оцењивање у логаритамском простору, повратна трансформација преко експоненцијалне
функције \\
\addlinespace
Ненегативни најмањи квадрати у нивоима &
\texttt{har} &
Ненегативни коефицијенти уз ненегативне регресоре \\
\addlinespace
Структурна мултипликативна форма &
\texttt{mem} &
Конвексна комбинација ненегативних чланова \\
\bottomrule
\end{tabularx}
\end{table}

\subsection{HAR модел и његове варијанте}

\paragraph{Полазна интуиција.}
Корси \parencite{corsi2009} полази од хипотезе о хетерогеном тржишту: учесници
реагују различитим брзинама (дневни трговци, недељни портфолио менаџери и дугорочни
институционални инвеститори), па волатилност представља каскаду компоненти различитих
временских хоризоната. Дуготрајно памћење тако се апроксимира без формалног
дугомеморијског процеса.

\paragraph{Логаритамски HAR.}
Стандардна форма и референтни модел целог рада:

\begin{equation}
\ln y_{t,h} = \beta_0 + \beta_d \ln\RV_{d,t} + \beta_w \ln\RV_{w,t} + \beta_m \ln\RV_{m,t} + \varepsilon_t
\label{eq:loghar}
\end{equation}

Оцењује се обичним методом најмањих квадрата, а прогноза се враћа у ниво преко
непараметарске корекције из одељка~\ref{subsec:jensen}. Овај модел, а не
\texttt{har}, јесте основа за релативни QLIKE.

\paragraph{HAR у нивоима са ненегативним најмањим квадратима.}
Иста скупина улазних променљивих, али потпуно другачији механизам позитивности:

\begin{equation}
y_{t,h} = \beta_0 + \beta_d \RV_{d,t} + \beta_w \RV_{w,t} + \beta_m \RV_{m,t} + \varepsilon_t
\end{equation}

уз ограничење да су сви коефицијенти ненегативни:

\begin{equation}
\hat{\beta} = \arg\min_{\beta \ge 0}\left\|y - X\beta\right\|^{2}
\end{equation}

Пошто су сви регресори ненегативни (реч је о варијансама) и сви коефицијенти
ненегативни, прогноза не може бити негативна по конструкцији. Нема логаритамске
трансформације ни корекције повратне трансформације. Ова два модела остају трајно
раздвојена: обједињавање у „један HAR“ избрисало би могућност поређења утицаја самог
механизма позитивности.

\paragraph{CHAR.}
Уместо $\RV$, као регресор користи се искључиво континуирана компонента:

\begin{equation}
\ln y_{t,h} = \beta_0 + \beta_d \ln C_{d,t} + \beta_w \ln C_{w,t} + \beta_m \ln C_{m,t} + \varepsilon_t
\end{equation}

\paragraph{HAR-J.}
Континуирана и скоковита компонента улазе као \textbf{одвојени} регресори, чиме се
тестира да ли скокови имају другачију перзистентност:

\begin{equation}
\ln y_{t,h} = \beta_0 + \beta_d \ln C_{d,t} + \beta_w \ln C_{w,t} + \beta_m \ln C_{m,t}
+ \beta_J \ln J_{d,t} + \varepsilon_t
\end{equation}

\paragraph{SHAR.}
Дневна компонента раздваја се по предзнаку, чиме се хвата асиметрија коју описују
Патон и Шепард:

\begin{equation}
\ln y_{t,h} = \beta_0 + \beta^{+}\ln\RS_{t}^{+} + \beta^{-}\ln\RS_{t}^{-}
+ \beta_w \ln\RV_{w,t} + \beta_m \ln\RV_{m,t} + \varepsilon_t
\end{equation}

\paragraph{HAR-IV.}
HAR проширен индексом имплицитне волатилности са доцњом од једног дана, чиме се
тестира да ли тржиште опција садржи информацију које нема у историји цена:

\begin{equation}
\ln y_{t,h} = \beta_0 + \beta_d \ln\RV_{d,t} + \beta_w \ln\RV_{w,t} + \beta_m \ln\RV_{m,t}
+ \beta_{IV}\VIX_{t-1} + \varepsilon_t
\end{equation}

\subsection{HARQ — модел изграђен на грешци мерења}
\label{subsec:harq}

HARQ \parencite{bpq2016} теоријски је најближи тези овог рада. Идеја је да коефицијент
уз дневну компоненту не буде константан, већ да зависи од тога колико је мерење тог
дана било поуздано:

\begin{equation}
\ln y_{t,h} = \beta_0 + \left(\beta_d + \beta_Q\sqrt{\RQ_t}\right)\ln\RV_{d,t}
+ \beta_w \ln\RV_{w,t} + \beta_m \ln\RV_{m,t} + \varepsilon_t
\label{eq:harq}
\end{equation}

Када је реализована квартичност висока, дакле када је мерење тог дана било зашумљено,
модел аутоматски смањује тежину те опсервације. Користи се квадратни корен реализоване
квартичности, а не сама вредност, како би интеракциони члан остао истог реда величине
као основни регресор.

Значај овог модела за рад је концептуалан: он представља постојећи, у литератури
признат модел изграђен управо на претпоставци да грешка мерења систематски утиче на
прогнозу. Овај рад ту идеју проширује тако што грешку мерења не моделује унутар једног
модела, већ је варира експериментално.

\subsection{ARFIMA — експлицитно дугорочно памћење}

За разлику од HAR-а, који дуготрајно памћење апроксимира каскадом, ARFIMA га моделује
експлицитно преко фракционог диференцирања:

\begin{equation}
(1-L)^{d}\left(\ln\RV_t - \mu\right) = \Psi(L)\varepsilon_t
\end{equation}

Оператор фракционог диференцирања развија се у биномни ред, чији се коефицијенти
рачунају рекурзивно:

\begin{equation}
(1-L)^{d} = \sum_{k=0}^{\infty}\binom{d}{k}(-L)^{k},
\qquad
w_k = w_{k-1}\cdot\frac{k-1-d}{k},
\quad w_0 = 1
\end{equation}

Параметар фракционог диференцирања $d$ оцењује се локалним Витловим естиматором
\parencite{robinson1995}, полупараметарским поступком који користи само ниске
фреквенције периодограма и стога не захтева претпоставку о целом спектралном облику:

\begin{equation}
\hat{d} = \arg\min_{d}\left[
\ln\!\left(\frac{1}{m}\sum_{j=1}^{m}\lambda_j^{2d}I(\lambda_j)\right)
- \frac{2d}{m}\sum_{j=1}^{m}\ln\lambda_j
\right]
\end{equation}

где је $I(\lambda_j)$ периодограм на фреквенцији $\lambda_j = 2\pi j/n$. Након
диференцирања, редови $(p,q)$ бирају се по Бајесовом информационом критеријуму преко
скупа $\{0,1,2\}^{2}$:

\begin{equation}
\BIC = -2\ln\hat{L} + k\ln n
\end{equation}

Прогноза се потом ре-интегрише у нивое инверзним оператором $(1-L)^{-d}$, чији се
тежински коефицијенти такође рачунају рекурзивно:

\begin{equation}
v_k = v_{k-1}\cdot\frac{k-1+d}{k},
\qquad v_0 = 1
\end{equation}

Важна имплементациона напомена: ре-интеграција мора бити извршена конволуцијом преко
целе доступне диференциране историје, а не као „последњи ниво плус локална промена“.
Оператор $(1-L)^{-d}$ је дугомеморијски, па његова реконструкција зависи од читаве
историје, и скраћивање би променило саму реконструисану вредност.

Модел се оцењује \textbf{независно по симболу}, за разлику од HAR породице која се
може обједињавати кроз попречни пресек.

\subsection{MEM — мултипликативни модел грешке}
\label{subsec:mem}

Енгл \parencite{engle2002} предлаже мултипликативну поставку у којој се реализована
варијанса представља као производ условне средине и позитивног случајног члана:

\begin{equation}
\RV_t = \mu_t\varepsilon_t,
\qquad \varepsilon_t > 0,
\qquad E[\varepsilon_t] = 1
\end{equation}

Условна средина прати рекурзију у GARCH стилу:

\begin{equation}
\mu_t = \omega + \alpha\RV_{t-1} + \beta\mu_{t-1}
\label{eq:mem}
\end{equation}

\textbf{Позитивност је структурна}: $\mu_t$ је конвексна комбинација ненегативних
чланова, па је $\mu_t \ge 0$ за свако $t$ по конструкцији, без икакве логаритамске
трансформације и без корекције повратне трансформације. То је трећи од три механизма
из одељка~\ref{subsec:pozitivnost}.

Оцењивање се врши квази-методом максималне веродостојности уз експоненцијалну
веродостојност:

\begin{equation}
\ell(\omega,\alpha,\beta) = \sum_{t}\left[-\ln\mu_t - \frac{\RV_t}{\mu_t}\right]
\end{equation}

уз ограничења $\omega > 0$, $\alpha \ge 0$, $\beta \ge 0$ и услов стационарности
$\alpha + \beta < 1$. Разлог за избор експоненцијалне веродостојности јесте робусност:
оцена остаје конзистентна за било коју стварну расподелу грешке са позитивним носачем
и јединичном средином, што је мултипликативни аналог Болерслевовог резултата за GARCH.

Прогноза за хоризонт $h$ добија се пројекцијом рекурзије унапред:

\begin{equation}
E\!\left[\mu_{t+1}\mid F_t\right] = \omega + \alpha\RV_t + \beta\mu_t
\end{equation}

\begin{equation}
E\!\left[\mu_{t+k}\mid F_t\right] = \omega + (\alpha+\beta)E\!\left[\mu_{t+k-1}\mid F_t\right],
\qquad k > 1
\end{equation}

а коначна прогноза јесте просек тих вредности преко $k = 1,\dots,h$, што одговара
конструкцији циља из одељка~\ref{subsec:target}.

\subsection{Модели засновани на стаблима}

\paragraph{Градијентно појачавање.}
LightGBM и XGBoost граде адитивни ансамбл плитких стабала одлучивања, где свако
наредно стабло учи на резидуалима претходних:

\begin{equation}
F_M(x) = \sum_{m=1}^{M}\nu f_m(x)
\end{equation}

где је $\nu$ стопа учења, а $f_m$ стабло оцењено на негативном градијенту функције
губитка. Циљна функција са регуларизацијом гласи

\begin{equation}
\mathcal{L} = \sum_{i}L\!\left(y_i, \hat{y}_i\right) + \sum_{m}\Omega(f_m),
\qquad
\Omega(f) = \gamma T + \frac{1}{2}\lambda\left\|w\right\|^{2}
\end{equation}

при чему $T$ означава број листова, а $w$ вредности у листовима.

Предност ових модела јесте способност хватања нелинеарних интеракција које линеарни
HAR не може представити. Улазни скуп представља унију свих променљивих HAR породице,
проширену полужним чланом:

\begin{equation}
X = \left\{\RV_d, \RV_w, \RV_m, C_d, C_w, C_m, J_d, \RS^{+}, \RS^{-},
\RV_d\sqrt{\RQ}, \VIX_{t-1}, r_d\right\}
\end{equation}

Хиперпараметри су \textbf{фиксирани и наведени}, а не подешавани по прозору (двеста
стабала, максимална дубина четири, петнаест листова, стопа учења $0{,}05$). Разлог је
методолошки: подешавани хиперпараметри учинили би да допринос варијанси одражава
уложен труд у претрагу, а не саму архитектуру, што би обесмислило поређење са осталим
факторима.

\subsection{PatchTST — трансформерска архитектура}

PatchTST \parencite{nie2023} прилагођава трансформерску архитектуру временским серијама
кроз два кључна елемента.

\paragraph{Сегментација на закрпе.}
Улазна секвенца дужине $L$ дели се на сегменте (\textit{patch}) дужине $P$, који
постају токени:

\begin{equation}
N_{patches} = \left\lfloor \frac{L}{P} \right\rfloor
\end{equation}

Сваки сегмент линеарно се пројектује у уграђени простор димензије $d_{model}$, уз
додавање позиционог уграђивања:

\begin{equation}
z_j = W_p x_j + e_j
\end{equation}

Овај поступак смањује број токена $P$ пута у односу на приступ у коме сваки временски
корак представља засебан токен, што квадратну сложеност механизма пажње чини
подношљивом.

\paragraph{Механизам пажње.}
Језгро трансформера \parencite{vaswani2017} јесте скалирана скаларна пажња:

\begin{equation}
\operatorname{Attention}(Q,K,V) = \operatorname{softmax}\!\left(\frac{QK^{\top}}{\sqrt{d_k}}\right)V
\end{equation}

где се упити, кључеви и вредности добијају линеарним пројекцијама уграђених сегмената.
Вишеглава пажња примењује овај поступак паралелно у више подпростора:

\begin{equation}
\operatorname{MultiHead}(Q,K,V) = \operatorname{Concat}(h_1,\dots,h_H)W^{O}
\end{equation}

\paragraph{Независност канала.}
Сваки улазни канал обрађује се независно, кроз исти енкодер са дељеним параметрима,
што спречава да модел научи лажне зависности између канала на узорку ограничене
величине.

Конфигурација у овом раду намерно је мала: дужина секвенце шездесет дана, дужина
сегмента десет, димензија модела тридесет два, четири главе пажње и два слоја
енкодера. Разлог је наведен у самом плану рада као мера ублажавања ризика:
трансформерске архитектуре тестиране су на скуповима са стотинама хиљада тачака, док
панел ове величине нуди знатно мање.

Улазни канали су $\{\RV_d, C_d, J_d, \RS^{+}, \RS^{-}, r_d\}$. Ово је једини
стохастички модел у раду, па се оцењује на \textbf{пет различитих семена}, уз
извештавање средње вредности и стандардне девијације. Без тога није могуће
разликовати ефекат архитектуре од шума иницијализације.

\subsection{Темељни модели временских серија}

Темељни (\textit{foundation}) модели представљају најновији правац: реч је о моделима
претходно обученим на великом броју разноврсних временских серија, који се потом
примењују \textbf{без икаквог додатног обучавања} (\textit{zero-shot}).

У овом раду користи се IBM Granite TinyTimeMixer, са \textbf{тачно прикованим}
контролним тачкама. Прикивање је методолошки неопходно: различита издања обучена су на
различитим скуповима података, па би провера контаминације над неприкованим моделом
била бесмислена. Реализоване варијанте приказане су у табели~\ref{tab:ttm}.

\begin{table}[htbp]
\centering
\caption{Варијанте темељног модела и њихова зависност од фактора мерења}
\label{tab:ttm}
\begin{tabular}{lll}
\toprule
Варијанта & Улаз & Зависи од \\
\midrule
\texttt{ttm\_rv}        & серија $\ln \RV_d$        & мреже \\
\texttt{ttm\_c}         & серија $\ln C_d$          & мреже, естиматора, теста \\
\texttt{ttm\_c\_plus\_j} & серије $C$ и $J$ одвојено & мреже, естиматора, теста \\
\texttt{log\_har\_ttm}  & хибрид са \texttt{log\_har} & мреже \\
\bottomrule
\end{tabular}
\end{table}

Модел \texttt{ttm\_c\_plus\_j} прогнозира континуирану компоненту на уобичајеном
логаритамском путу, а скоковиту компоненту \textbf{у нивоима}, јер $J_t = 0$ важи на
већини дана, па је $\ln J$ недефинисан. Компоненте се сабирају у нивоима и на збир се
примењује \textbf{једна} комбинована корекција из одељка~\ref{subsec:jensen}.

Хибридни модел комбинује прогнозе у логаритамском простору, пре корекције, са тежинама
\textbf{фиксираним на $0{,}5$ и $0{,}5$ и никада оптимизованим}:

\begin{equation}
\widehat{\ln y}^{\,hybrid} = 0{,}5\,\widehat{\ln y}^{\,log\_har} + 0{,}5\,\widehat{\ln y}^{\,ttm}
\end{equation}

Разлог за фиксиране тежине идентичан је разлогу за фиксиране хиперпараметре стабала:
оптимизована комбинација одражавала би уложен труд у претрагу, а не архитектуру.

\textbf{Провера контаминације предтренирања} представља првокласни, забележени
артефакт. Она пореди објављени датум одсецања података за предтренирање са тест
периодом и уписује резултат \textbf{у оба случаја}, чист или не. Када се контаминација
утврди, преклапање се квантификује, а модел се \textbf{задржава у анализи уз ознаку да
је контаминиран}: никада се тихо не избацује нити тихо третира као чист.

\section{Валидациони протокол}

\subsection{Усидрена унапредна валидација}
\label{subsec:validacija}

Свака подела података мора искључити свако цурење информација из будућности. Користи
се \textbf{усидрена} (\textit{anchored}) унапредна валидација, код које је почетак
тренинга фиксиран, а само се крај тренинга шири:

\begin{equation}
\text{прозор } w:\quad
\text{тренинг} = \left[T_0,\ T_0 + 5 + w\right],
\qquad
\text{тест} = \left[T_0 + 5 + w,\ T_0 + 6 + w\right]
\end{equation}

Укупно се формира једанаест прозора над периодом 2010--2025: први прозор тренира на
периоду 2010--2014 и тестира на 2015, а једанаести тренира на 2010--2024 и тестира на
2025. Валидација никада није клизна, јер фиксирано сидро обезбеђује да се сав
расположиви историјски узорак увек користи.

\subsection{Валидациони преклоп}

Унутар сваког прозора, последња година тренинга издваја се као \textbf{валидациони
преклоп}. Он се користи искључиво за две сврхе:

\begin{enumerate}
\item оцену непараметарске корекције повратне трансформације
      (одељак~\ref{subsec:jensen}),
\item оцену Минцер-Зарновиц коефицијената за рекалибрацију
      (одељак~\ref{subsec:mz}).
\end{enumerate}

Никада се не користи за подешавање хиперпараметара нити за евалуацију. Ова
раздвојеност је суштинска: коришћење истих података за оцену корекције и за
извештавање резултата представљало би прикривени облик гледања унапред.

\subsection{Обрада стохастичких модела}

За PatchTST, једини стохастички модел, примарна извештавана величина усредњава
\textbf{губитке по семену}, а не прогнозе:

\begin{equation}
\bar{L} = \frac{1}{S}\sum_{s=1}^{S}L\!\left(y, \hat{y}^{(s)}\right)
\end{equation}

Ансамбл усредњених прогноза постоји као експлицитно \textbf{одвојена, секундарна}
величина и никада не замењује примарну. Код ансамбла, усредњавање се врши у
логаритамском простору, пре повратне трансформације:

\begin{equation}
\widehat{\ln y}^{\,ens} = \frac{1}{S}\sum_{s=1}^{S}\widehat{\ln y}^{\,(s)},
\qquad
\bar{a} = \frac{1}{S}\sum_{s=1}^{S}a^{(s)}
\end{equation}

\begin{equation}
\hat{y}^{\,ens} = \exp\!\left(\widehat{\ln y}^{\,ens} + \bar{a}\right)
\end{equation}

Усредњавање већ трансформисаних прогноза двоструко би урачунало корекцију сваког
семена.

\section{Евалуација}

Метрике су организоване у четири нивоа, јер мере суштински различите ствари.

\subsection{Ниво 1 — тачност прогнозе}

\paragraph{Фундаментални проблем.}
Стварна волатилност је немерљива, па се прогнозе пореде у односу на зашумљен прокси.
Патон \parencite{patton2011} показао је да већина уобичајених функција губитка даје
\textbf{погрешно рангирање модела} када је прокси зашумљен, и извео потребне и довољне
услове које функција губитка мора испунити да би рангирање било робусно.

\paragraph{QLIKE (примарна метрика).}
Робусна је на шум у проксију:

\begin{equation}
\QLIKE_t = \frac{y_t}{\hat{y}_t} - \ln\!\left(\frac{y_t}{\hat{y}_t}\right) - 1
\label{eq:qlike}
\end{equation}

Губитак је нула само када се прогноза тачно поклопи са оствареном вредношћу. Функција
је \textbf{асиметрична}: потцењивање волатилности кажњава се јаче од прецењивања, што
је економски исправно, јер је потценити ризик скупље.

\paragraph{RMSE (секундарна).}
Такође робусна на шум у проксију, али симетрична:

\begin{equation}
\RMSE = \sqrt{\frac{1}{N}\sum_{t=1}^{N}\left(y_t - \hat{y}_t\right)^{2}}
\end{equation}

\paragraph{MAE (додатна).}
Није робусна на шум у проксију:

\begin{equation}
\MAE = \frac{1}{N}\sum_{t=1}^{N}\left|y_t - \hat{y}_t\right|
\end{equation}

Важна методолошка напомена: QLIKE и RMSE минимизовани су \textbf{условном средином}, а
MAE \textbf{условном медијаном}. Поређење модела по MAE над средински-коригованим
прогнозама кажњавало би модел зато што ради тачно оно што QLIKE захтева. Зато се као
обавезан партнер рачуна варијанта са корекцијом $a = 0$:

\begin{equation}
\MAE^{med} = \frac{1}{N}\sum_{t=1}^{N}\left|y_t - \exp\!\left(\widehat{\ln y}_t\right)\right|
\end{equation}

\paragraph{Релативни QLIKE.}
Приказна мера која резултате чини одмах читљивим:

\begin{equation}
\text{relQLIKE} = \frac{\overline{\QLIKE}_{model}}{\overline{\QLIKE}_{log\_har}}
\label{eq:relqlike}
\end{equation}

Ово је \textbf{однос средина, никада средина односа}. Разлика није техничка ситница:
средина односа експлодира чим губитак референтног модела на једном јединном дану приђе
нули, што се дешава, па један дан може преузети цео извештани број. Референтни модел је
\texttt{log\_har}, а вредност $0{,}95$ значи 5\% боље од референтног.

\subsection{Ниво 2 — статистичка значајност}

\paragraph{Дибold-Маријанов тест.}
За два модела рачуна се дневна разлика губитака \parencite{dieboldMariano1995}:

\begin{equation}
d_t = L_t^{(A)} - L_t^{(B)}
\end{equation}

Тест статистика пореди просек те разлике са њеном дугорочном стандардном грешком:

\begin{equation}
\DM = \frac{\bar{d}}{\sqrt{\hat{V}(\bar{d})/N}}
\end{equation}

Пошто су циљне вредности на $h = 5$ и $h = 22$ преклапајући просеци, разлика губитака
је аутокорелисана, па се користи Њуи-Вестова HAC варијанса \parencite{neweyWest1987} са
$h-1$ померајем:

\begin{equation}
\hat{V}(\bar{d}) = \gamma_0 + 2\sum_{k=1}^{h-1}\left(1 - \frac{k}{h}\right)\gamma_k
\end{equation}

Примењује се и Харви-Лејборн-Нјуболдова корекција за мале узорке
\parencite{harvey1997}:

\begin{equation}
\DM^{*} = \DM\sqrt{\frac{N + 1 - 2h + h(h-1)/N}{N}}
\end{equation}

\paragraph{Контрола вишеструког тестирања.}
Са тридесет шест дизајна и око четрнаест модела, исцрпно поређење свих парова даје
хиљаде тестова, па би се при $\alpha = 0{,}05$ стотине „значајних“ резултата појавило
чисто случајно. Свака извештавана \textbf{породица} тестова стога пролази кроз
Бенјамини-Хохбергову контролу стопе лажних открића, истом процедуром описаном у
одељку~\ref{subsec:a3}. Породица се мора навести експлицитно и бележи се уз резултат,
јер контрола стопе лажних открића нема смисла без изјаве над чим се контролише.

\paragraph{Скуп поверења модела.}
Хансен, Лунде и Насон \parencite{hansen2011} предлажу поступак који полази од свих
модела и итеративно одбацује најгори док се преостали не могу статистички разликовати
међусобно.

Суштинска напомена, коју рад мора експлицитно навести, јесте да \textbf{припадност
скупу поверења није рангирање тачкастих губитака}. Модел са вишим просечним QLIKE-ом
од неког елиминисаног модела може ипак остати у скупу, јер елиминација зависи од односа
разлике губитака према њеној стандардној грешци, а не од саме разлике: стабилан модел
који је благо лошији преживљава, док модел са бољим просеком али нестабилном серијом
разлика бива одбачен. Унутар скупа не постоји уређење, а величина скупа мери
\textbf{снагу узорка}, а не квалитет модела.

\subsection{Ниво 2 — декомпозиција варијансе}
\label{subsec:anova}

Ово је централни резултат рада и одељак који даје бројчан одговор на питање из наслова.

\textbf{Зависна променљива} јесте \textbf{дневни} QLIKE губитак, један ред по
комбинацији (модел, дизајн, хоризонт, симбол, дан). Коришћење средина ћелија не би
оставило никакав члан грешке, па оцена не би била могућа.

\textbf{Модел декомпозиције} за моделе осетљиве на дизајн гласи:

\begin{equation}
\begin{split}
L_{ijklmn} = {}& \mu + \tau_i^{model} + \gamma_j^{grid} + \delta_k^{est}
+ \psi_l^{test} + \rho_m^{regime} + \phi_n^{symbol} \\
& + (\gamma\delta)_{jk} + (\gamma\psi)_{jl} + (\tau\psi)_{il} + \varepsilon_{ijklmn}
\end{split}
\label{eq:anova}
\end{equation}

Четири својства ове поставке методолошки су суштинска.

\paragraph{1. Оцењивање по хоризонту, никада обједињено.}
Хоризонт се налази у кључу реда, али \textbf{није} фактор у формули. Разлог је то што
$h = 1$, $h = 5$ и $h = 22$ имају различите скале циљне величине, па би обједињавање
бацило објашњиву варијансу у члан грешке и умањило сваки парцијални ета-квадрат,
укључујући управо удео мерења ради којег рад постоји.

\paragraph{2. Симбол као експлицитан блокирајући фактор.}
Остављен у резидуалу, симбол би надувао члан грешке и смањио све извештаване уделе.

\paragraph{3. Сума квадрата типа III са сум-то-зеро контрастима.}
Дизајн је неуравнотежен по конструкцији, јер модели који зависе само од мреже добијају
дуплиране редове преко дванаест комбинација естиматора и теста, и садржи три намерна
интеракциона члана. При таквом дизајну тип II и тип III не слажу се око главних
ефеката. Тип III без сум-то-зеро кодирања није добро дефинисан, јер главни ефекти
зависе од произвољног избора референтне категорије, па се контрастно кодирање поставља
експлицитно:

\begin{equation}
\sum_{i}\tau_i = 0,
\qquad
\sum_{j}\gamma_j = 0,
\qquad \text{и тако редом}
\end{equation}

\paragraph{4. Интеракциони чланови су структурни, а не случајни.}
Микроструктурни шум подиже $\RV$ приближно линеарно са бројем интрадневних опсервација
(једначина~\ref{eq:noise-bias}), а естиматор отпоран на скокове погођен је тим истим
шумом другачије, па понашање естиматора у односу на $\RV$ заиста зависи од мреже.
Изостављање тих чланова тихо би пребацило део ефекта у главне ефекте.

\textbf{Мера удела варијансе} јесте парцијални ета-квадрат:

\begin{equation}
\eta_p^{2} = \frac{SS_{effect}}{SS_{effect} + SS_{error}}
\label{eq:eta}
\end{equation}

\paragraph{Ограничење које мора бити наведено у тексту рада.}
Различити дизајни мерења изведени су из \textbf{истих} сирових једноминутних свећица,
па су вредности $\RV$ на мрежама $1m$ и $5m$ за исти дан високо корелисане.
Претпоставка независности стандардног $F$-теста овде \textbf{не важи}. Последице су
следеће:

\begin{itemize}
\item Парцијални ета-квадрат извештава се као \textbf{дескриптиван} удео варијансе,
      што не захтева независност.
\item Сваки извештани $F$ или $p$ мора се читати уз \textbf{двосмерно кластерисане}
      стандардне грешке \parencite{cgm2011}, кластерисане истовремено по дану и по
      симболу:
      \begin{equation}
      \hat{V}_{2way} = \hat{V}_{day} + \hat{V}_{symbol} - \hat{V}_{day \cap symbol}
      \end{equation}
\item Као алтернативна инференцијална путања користи се блок бутстрап преко дана, са
      дужином блока \textbf{једнаком хоризонту} $h$, истом конвенцијом коју користи и
      скуп поверења модела, како се две путање не би разилазиле по конструкцији.
\end{itemize}

\paragraph{Робусност на асиметрију је обавезна, а не опциона.}
Дневни QLIKE изразито је десно асиметричан, а узорак обухвата 2018. годину и период
пандемије, па шачица дана може одредити уделе варијансе док излазна табела изгледа
сасвим уредно. Иста декомпозиција стога се изводи \textbf{четири начина, један поред
другог}, како је приказано у табели~\ref{tab:robusnost}.

\begin{table}[htbp]
\centering
\caption{Варијанте провере робусности декомпозиције варијансе}
\label{tab:robusnost}
\begin{tabularx}{\textwidth}{L{3.6cm} L{3.8cm} X}
\toprule
Варијанта & Зависна променљива & Сврха \\
\midrule
\texttt{raw}             & дневни QLIKE          & примарна \\
\texttt{log}             & $\ln(\QLIKE)$          & смањује утицај репова \\
\texttt{within\_day\_rank} & ранг унутар дана     & потпуно непараметарска \\
\texttt{raw\_ex\_top1pct} & без горњег 1\% дана   & одговор на питање да ли налаз важи
и без марта 2020. године \\
\bottomrule
\end{tabularx}
\end{table}

Стабилни удели кроз све четири варијанте представљају јаку тврдњу о робусности, док
нестабилност мора бити видљива \textbf{пре} него што се текст напише.

\subsection{Минцер-Зарновиц регресија}
\label{subsec:mz}

Регресија остварене вредности на прогнозу \parencite{mincerZarnowitz1969} открива
систематску пристрасност коју QLIKE не издваја посебно:

\begin{equation}
y_t = \alpha + \beta\hat{y}_t + u_t
\end{equation}

Тестирају се хипотезе $\alpha = 0$ и $\beta = 1$. Модел може имати низак губитак, а
доследно потцењивати ниво волатилности. Пошто су резидуали на дужим хоризонтима
аутокорелисани по конструкцији, користе се Њуи-Вестове HAC стандардне грешке са $h-1$
померајем.

Коефицијенти се оцењују \textbf{искључиво на валидационом преклопу}, а рекалибрисани
QLIKE извештава се \textbf{упоредо} са сировим, никада уместо њега, јер би
рекалибрација пре примарне метрике представљала гледање унапред.

\subsection{Ниво 3 — економска евалуација}

Прогноза волатилности користи се за израчунавање вредности при ризику:

\begin{equation}
\VaR_t(\alpha) = z_{1-\alpha}\sqrt{\hat{y}_t}
\end{equation}

где је $z_{1-\alpha}$ квантил стандардне нормалне расподеле ($2{,}33$ за ниво од 1\%).
Користи се квадратни корен прогнозиране варијансе, дакле стандардна девијација.

\paragraph{Кључно методолошко усклађивање.}
Пошто се $\RV$ рачуна искључиво из интрадневних приноса, а прекноћни принос је намерно
искључен, $\sqrt{\hat{y}}$ представља стандардну девијацију \textbf{само интрадневног}
кретања. Пробијања се стога морају мерити у односу на принос од отварања до затварања,
никада у односу на серију од затварања до затварања. Коришћење ове друге систематски
би потценило ризик, па би Купијецов тест одбацивао хипотезу због неусклађености
јединица, а не због стварног налаза.

Индикатор пробијања дефинише се као

\begin{equation}
I_t = \mathbf{1}\left\{-r_{d,t} > \VaR_t(\alpha)\right\}
\end{equation}

\paragraph{Купијецов тест безусловне покривености.}
Проверава да ли број пробијања одговара декларисаном нивоу \parencite{kupiec1995}:

\begin{equation}
LR_{uc} = -2\ln\!\left[\frac{\alpha^{x}(1-\alpha)^{N-x}}
{\hat{\pi}^{x}(1-\hat{\pi})^{N-x}}\right]\ \sim\ \chi^{2}(1)
\end{equation}

где је $x$ број пробијања, а $\hat{\pi} = x/N$ њихов емпиријски удео.

\paragraph{Кристоферсенов тест независности.}
Проверава да ли вероватноћа пробијања зависи од тога да ли је претходни дан био
пробијање \parencite{christoffersen1998}. Ако се са $n_{ij}$ означи број прелаза из
стања $i$ у стање $j$:

\begin{equation}
LR_{ind} = -2\ln\!\left[
\frac{\hat{\pi}^{\,n_{01}+n_{11}}(1-\hat{\pi})^{\,n_{00}+n_{10}}}
{\hat{\pi}_{01}^{\,n_{01}}(1-\hat{\pi}_{01})^{\,n_{00}}
\hat{\pi}_{11}^{\,n_{11}}(1-\hat{\pi}_{11})^{\,n_{10}}}
\right]\ \sim\ \chi^{2}(1)
\end{equation}

Концептуални значај овог теста је велик: модел може имати тачно очекиван број
пробијања, али ако се сва догоде исте недеље, реч је о катастрофалном моделу ризика.

\paragraph{Динамички квантилни тест.}
Представља строжу верзију провере независности \parencite{engleManganelli2004}.
Величина „погодак“ регресује се на сопствене доцње и на сам ниво вредности при ризику,
па се тестира да ли су сви коефицијенти једнаки нули:

\begin{equation}
\Hit_t = I_t - \alpha,
\qquad
\Hit_t = \lambda_0 + \sum_{k=1}^{K}\lambda_k \Hit_{t-k} + \lambda_{K+1}\VaR_t + u_t
\end{equation}

\paragraph{Лопезова функција губитка.}
Пробијања не само да броји, већ их кажњава сразмерно величини премашаја
\parencite{lopez1999}:

\begin{equation}
L^{Lopez} = \frac{1}{N}\sum_{t=1}^{N}
\left[1 + \left(\text{губитак}_t - \VaR_t\right)^{2}\right]I_t
\end{equation}

Ово разликује модел који греши мало од модела који греши катастрофално, дакле разлику
коју сам број пробијања не види.

Тестови се рачунају на \textbf{обједињеној унији свих тест прозора} (приближно 2.772
дана кроз једанаест прозора), а не по појединачном прозору: један прозор од 252 дана
очекује свега око $2{,}5$ пробијања на нивоу од 1\%, што Купијецовом тесту даје
суштински никакву снагу.

\subsection{Ниво 4 — карактеризација детекције скокова}

\paragraph{Методолошка напомена.}
Класификационе метрике (прецизност, одзив, $F_1$) овде \textbf{нису применљиве}, јер не
постоји основна истина: скок није опсервабилан догађај, него оно што естиматор прогласи
скоком. Уместо тачности, мере се слагање и последице.

\paragraph{Жакаров индекс преклапања.}
Мери колико се две варијанте теста међусобно слажу:

\begin{equation}
J(A,B) = \frac{\left|A \cap B\right|}{\left|A \cup B\right|}
\end{equation}

Низак индекс значи да избор теста суштински мења улазне податке, што је сама премиса
рада.

\paragraph{Удео дана са скоком.}
Мери колико је метода конзервативна:

\begin{equation}
\text{удео} = \frac{\text{број дана са детектованим скоком}}{\text{укупан број дана}}
\end{equation}

Литература наводи типичан опсег од 5\% до 15\%, па вредност изван тог опсега
сигнализира лоше подешен праг.

\paragraph{Удео скоковите компоненте у укупној варијанси.}
Мери економску значајност:

\begin{equation}
\text{удео } J = \frac{\sum_t J_t}{\sum_t \RV_t}
\end{equation}

\paragraph{Плацебо тест.}
Индикатори скокова насумично се пермутују кроз време, чиме се задржава њихова
учесталост, али уништава веза са стварним данима. Ако побољшање опстане и са
пермутованим индикаторима, оно \textbf{није} последица информације о скоковима, него
артефакт структуре модела.

Важно ограничење јесте да је QLIKE недефинисан при $J_t = 0$ (дељење нулом и логаритам
нуле), па се никада не примењује на серију скоковите компоненте, већ искључиво на
прогнозу реализоване варијансе.

\subsection{Стратификација резултата}
\label{subsec:stratifikacija}

Просек кроз цео узорак може сакрити да модел побеђује у мирним, а губи у турбулентним
периодима, што је за управљање ризиком управо обрнуто од корисног. Све метрике рачунају
се одвојено по три димензије, приказане у табели~\ref{tab:stratifikacija}.

\begin{table}[htbp]
\centering
\caption{Димензије стратификације резултата}
\label{tab:stratifikacija}
\begin{tabular}{ll}
\toprule
Димензија & Подела \\
\midrule
Режим            & мирни наспрам турбулентних периода, по квантилу $\RV_m$ \\
Присуство скока  & дани са скоком наспрам дана без скока \\
Хоризонт         & $h = 1$, $h = 5$, $h = 22$ \\
\bottomrule
\end{tabular}
\end{table}

Граница режима рачуна се \textbf{једном}, искључиво из тренинг података првог прозора,
и трајно се памти. Она се никада не рачуна изнова из података прозора који је тренутно
у обради: када би се граница померала по прозору, ознака режима за дати календарски дан
постала би двосмислена, јер би зависила од тога који је прозор произвео прогнозу која
се означава.

\section{Празнина у постојећој литератури}
\label{sec:praznina}

Преглед области показује да постоје три групе радова, од којих ниједна не одговара на
питање које овај рад поставља.

Прву групу чине радови који уводе нове економетријске моделе, најчешће проширења
хетерогеног ауторегресивног модела. Они пореде своје проширење са постојећим
варијантама, али све моделе оцењују на јединственом, унапред изабраном скупу
реализованих мера.

Другу групу чине радови који примењују методе машинског и дубоког учења. Они проширују
скуп архитектура, али задржавају исту прећутну претпоставку: улазне мере су фиксне, а
вари́ра се само модел.

Трећу групу чине радови из области мерења, који уводе нове естиматоре или тестове за
детекцију скокова. Они пажљиво испитују својства саме мере, али последице избора мере
по тачност предвиђања обично не прате низводно, кроз више модела и хоризоната.

Изузетак у односу на прву групу представља модел HARQ, описан у
одељку~\ref{subsec:harq}. Он изричито узима у обзир чињеницу да мерење садржи грешку и
да та грешка варира кроз време, чиме потврђује да премиса овог рада није нова ни
ексцентрична. Разлика је у томе што HARQ грешку мерења моделује унутар једног модела,
уместо да је третира као димензију која се може експериментално варирати и чији се
допринос може измерити.

Празнина је, дакле, следећа: не постоји систематско поређење у коме се дизајн мерења
третира као експериментална оса равноправна избору модела, а њихови доприноси тачности
раздвајају бројчано.

\section{Синтеза теоријског оквира}

Претходни одељци успостављају два одвојена извора варијације у тачности прогнозе.

Са једне стране, \textbf{дизајн мерења} чине три независне одлуке (фреквенција
узорковања, естиматор отпоран на скокове и тест за детекцију скокова), које дају
тридесет шест различитих, а подједнако легитимних начина да се из истих сирових
података конструишу исте номиналне улазне величине. Свака од тих одлука мења вредности
$\RV$, $C$ и $J$ које модел види.

Са друге стране, \textbf{избор модела} обухвата распон од наивне основе, преко
економетријске HAR породице и њених проширења, дугомеморијских и мултипликативних
модела, до модела заснованих на стаблима, трансформерске архитектуре и темељних модела
примењених без обучавања.

Постојећа литература варира искључиво другу димензију, док прву третира као фиксну и
унапред дату. HARQ модел из одељка~\ref{subsec:harq} представља важан изузетак, јер
грешку мерења узима у обзир, али је моделује унутар једног модела.

Овај рад обе димензије третира равноправно и укршта их у потпуни факторски експеримент.
Декомпозиција варијансе из одељка~\ref{subsec:anova} представља апарат који разлике у
губитку приписује појединачним факторима, чиме питање да ли мерење утиче више од
архитектуре престаје да буде тврдња и постаје број који се може навести и проверити.

Преостала два елемента теоријског оквира не улазе у то укрштање, али су за
тумачење резултата неопходна. Изведене променљиве из одељка~\ref{sec:izvedene}
посредују између двају фактора: оне настају из дизајна мерења и представљају
једини облик у ком га модели уопште виде, због чега промена дизајна и може да
се одрази на тачност. Стратификација из одељка~\ref{subsec:stratifikacija}
дели већ добијене резултате по режиму, присуству скока и хоризонту, чиме
открива да ли утврђени удели варијансе важе равномерно или само у појединим
тржишним условима. Заједно, ова четири елемента чине оквир унутар кога се
централно питање рада може не само поставити, него и бројчано проверити.
